% Week 4: Analysis and Reporting – Empirical Strategies for Data Interpretation
% BA-Eindwerkstuk Seminar - Leiden University
% Dr. Steven Denney

\documentclass[aspectratio=169,11pt]{beamer}

% ============================================
% Packages
% ============================================
\usepackage{tikz}
\usepackage{graphicx}
\usepackage{hyperref}
\usepackage{fontspec}
\usepackage{booktabs}

% Load custom theme
\usepackage{beamerthemethesis}

% ============================================
% Metadata
% ============================================
\title{BA-Eindwerkstuk Seminar}
\subtitle{Week 4: Analysis and Reporting}
\author{Dr. Steven Denney}
\institute{Korean Studies \\ Leiden University}
\date{February 27, 2026}

% ============================================
% Document
% ============================================
\begin{document}

% --------------------------------------------
% Title Slide
% --------------------------------------------
\begin{frame}[plain,noframenumbering]
    \titlepage
\end{frame}

% ============================================
\section{Today's Session}
% ============================================

\begin{frame}{How Today Works}
    Dr.\ Denney is out sick. Today's session is led by your classmates.
    \vspace{0.5cm}

    \textbf{The plan:}
    \begin{itemize}
        \item As a class, go through these slides together and talk through them as a group.
        \item No one is lecturing --- we are reading, reacting, and discussing.
        \item Please raise questions, comments, and concerns as we go.
    \end{itemize}
    \vspace{0.3cm}
    \begin{alertblock}{Important}
        We are taking notes on what gets discussed, asked, and raised. Dr.\ Denney will respond in writing or by video afterward. Your questions and comments will help shape that response.
    \end{alertblock}
\end{frame}

\begin{frame}{What You Need}
    Download these two documents from the course website \textbf{now}:
    \vspace{0.5cm}
    \begin{enumerate}
        \item \textbf{Analytical Framework Guidelines} --- explains the structure you will use in your thesis
        \vspace{0.3cm}
        \item \textbf{Analytical Framework Example} --- an example showing how a real project uses this structure
    \end{enumerate}
    \vspace{0.5cm}
    We will read the \textbf{guidelines first}, then come back to the example later.
\end{frame}

% ============================================
\section{The Analytical Framework}
% ============================================

\begin{frame}{What Is an Analytical Framework?}
    An analytical framework is a \textbf{structured plan} that explains how you will interpret your data and answer your research question.
    \vspace{0.3cm}

    It connects three things you have already developed:
    \vspace{0.3cm}
    \begin{enumerate}
        \item Your \textbf{research question} (Week 1)
        \item Your \textbf{data and sources} (Week 2)
        \item Your \textbf{literature review} and theoretical approach (Week 3)
    \end{enumerate}
    \vspace{0.3cm}
    \begin{exampleblock}{A note on terminology}
        Depending on your project, this chapter might be called ``Methodology,'' ``Research Design,'' or ``Analytical Approach.'' ``Analytical framework'' is a generic term --- what matters is the content, not the label.
    \end{exampleblock}
\end{frame}

\begin{frame}{Where It Fits in Your Thesis}
    \begin{enumerate}
        \item \textbf{Introduction} --- question, motivation, overview
        \item \textbf{Literature Review} --- situate your question, identify the gap
        \item \textbf{Methodology / Analytical Framework} $\leftarrow$ \textit{this is what today is about}
        \begin{itemize}
            \item How you connect theory to data
            \item What you analyze, how, and why
        \end{itemize}
        \item \textbf{Findings}
        \item \textbf{Discussion and Conclusion}
    \end{enumerate}
    \vspace{0.3cm}
    The analytical framework is the bridge between your literature review and your findings. Without it, the reader doesn't know how you got from theory to evidence.
\end{frame}

\begin{frame}[plain,noframenumbering]
    \vspace{2.5cm}
    \begin{center}
        {\color{TextMuted}\large Read}
        \vspace{0.2cm}

        {\color{AccentBlue}\Large\bfseries Analytical Framework Guidelines}
        \vspace{0.3cm}

        {\color{DustyRose}\rule{3cm}{0.8pt}}
        \vspace{0.5cm}

        {\color{TextMuted}\normalsize This document walks through a suggested structure for writing\\the methodology chapter of your thesis. Take a few minutes to read it now.}
    \end{center}
\end{frame}

\begin{frame}{Discuss: The Guidelines}
    After reading through the guidelines document:
    \vspace{0.3cm}
    \begin{itemize}
        \item What parts of this structure make sense to you? What parts are confusing or unclear?
        \vspace{0.3cm}
        \item Which sections do you think will be easiest to write for your project? Which will be hardest?
        \vspace{0.3cm}
        \item The guidelines note that not every section will look the same for every project. For example, ``operationalization'' applies most directly to projects with coding schemes --- if your project is more historical or interview-based, the equivalent is explaining how you identify themes or patterns in your sources. What does this look like for your project?
    \end{itemize}
    \vspace{0.3cm}
    \textit{Note any questions you want Dr.\ Denney to address.}
\end{frame}

% ============================================
\section{Seeing It in Practice}
% ============================================

\begin{frame}[plain,noframenumbering]
    \vspace{2.5cm}
    \begin{center}
        {\color{TextMuted}\large Read}
        \vspace{0.2cm}

        {\color{AccentBlue}\Large\bfseries Analytical Framework Example}
        \vspace{0.3cm}

        {\color{DustyRose}\rule{3cm}{0.8pt}}
        \vspace{0.5cm}

        {\color{TextMuted}\normalsize This document applies the guidelines structure to a real project ---\\media coverage of Yemeni refugees on Jeju Island. Take a few minutes to read it now.}
    \end{center}
\end{frame}

\begin{frame}{What Makes This Example Work?}
    \begin{itemize}
        \item \textbf{Theory drives the method.} The four coding categories (security, economic burden, cultural identity, humanitarian) come from framing theory --- they are not made up. The researcher can explain \textit{why} those categories.
        \vspace{0.3cm}
        \item \textbf{Sources are justified.} Three newspapers chosen because they represent different ideological positions. The comparison \textit{is} the analysis.
        \vspace{0.3cm}
        \item \textbf{Steps are explicit.} Someone else could follow the same procedure and arrive at similar findings.
        \vspace{0.3cm}
        \item \textbf{Everything connects back to the question.} Every element traces back to the research question about how ideological stances shape media representation.
    \end{itemize}
\end{frame}

\begin{frame}{Discuss: The Example}
    \vspace{0.3cm}
    \begin{itemize}
        \item Does this framework make sense? What would you do differently?
        \vspace{0.3cm}
        \item The coding scheme has four categories. How did the researcher get from ``framing theory'' to \textit{these specific four}? Is that convincing?
        \vspace{0.3cm}
        \item Think about your own project: what is the equivalent of ``comparing three newspapers''? What are you comparing --- sources, time periods, cases, perspectives?
        \vspace{0.3cm}
        \item What questions do you have for Dr.\ Denney about the example or about building your own framework?
    \end{itemize}
\end{frame}

% ============================================
\section{Your Projects}
% ============================================

\begin{frame}{Thinking About Your Own Framework}
    Try to answer these four questions for your project --- even roughly:
    \vspace{0.3cm}
    \begin{enumerate}
        \item \textbf{What is your theoretical approach?} What theory or concept from your literature will guide your analysis?
        \vspace{0.2cm}
        \item \textbf{What are your sources?} What will you actually analyze, and why those?
        \vspace{0.2cm}
        \item \textbf{How will you connect theory to data?} How will you recognize your key concepts in the data? What will you look for?
        \vspace{0.2cm}
        \item \textbf{What are your steps?} What do you do first, second, third?
    \end{enumerate}
    \vspace{0.3cm}
    Go around the group. Share your answers --- even if they are rough or uncertain. That is the point. Note where people are stuck or unsure.
\end{frame}

% ============================================
\section{Next Steps}
% ============================================

\begin{frame}{Assignment \#1: Revised Proposal}
    \textbf{Due: March 13, 2026 --- submit via Brightspace by 23:59.}
    \vspace{0.3cm}

    Your revised proposal should incorporate what you have learned in weeks 1--4:
    \vspace{0.3cm}
    \begin{itemize}
        \item A clear, focused \textbf{research question} and statement of the \textbf{problem}
        \item An \textbf{annotated bibliography} and draft \textbf{literature review outline}
        \item A description of your \textbf{data and sources}
        \item A preliminary \textbf{analytical framework} (use the guidelines as a reference)
    \end{itemize}
    \vspace{0.3cm}
    \begin{block}{Rubric}
        See the rubric on the course website for full evaluation criteria.
    \end{block}
\end{frame}

\begin{frame}{Looking Ahead}
    \textbf{This is the last weekly ``boot camp'' session.}
    \vspace{0.3cm}

    \begin{itemize}
        \item \textbf{Week 6 (Mar.\ 20):} Pre-recorded lecture reviewing weeks 1--4
        \item \textbf{From now on:} Work independently; continue to meet with your supervisors
    \end{itemize}
    \vspace{0.3cm}
    \begin{block}{Office Hours --- Next Week}
        Dr.\ Denney will be in the classroom \textbf{Friday, March 6, from 10:00 to 11:00} for questions and answers --- about this week's material, your proposals, or anything else.
    \end{block}
    \vspace{0.3cm}
    \textbf{Key Dates:}
    \begin{itemize}
        \item Assignment \#1 due: March 13, 2026
        \item Assignment \#2 due: April 03, 2026
    \end{itemize}
\end{frame}

\begin{frame}{Resources}
    Analytical Framework Guidelines (on the course website)\\
    \vspace{0.3cm}

    Analytical Framework Example: Yemeni Refugees (on the course website)\\
    \vspace{0.5cm}

    Mullaney, Thomas S., and Christopher Rea. \textit{Where Research Begins: Choosing a Research Project That Matters to You (and the World)}. Chicago: University of Chicago Press, 2022.\\
    \vspace{0.5cm}

    Assignment \#1 rubric available on the course website.
\end{frame}

\end{document}
